% =============================================================================
% SECTION 5: THE LOSCHMIDT PARADOX AS CATEGORY ERROR
% =============================================================================

\section{The Loschmidt Paradox as a Category Error}
\label{sec:loschmidt}

\subsection{Statement of the Paradox}
\label{subsec:paradox_statement}

Loschmidt's paradox arises from the following argument. The microscopic equations
of motion of classical and quantum mechanics are time-reversal symmetric: if
$\{x_i(t), p_i(t)\}$ is a solution, then $\{x_i(-t), -p_i(-t)\}$ is also a
solution. Therefore, if a system evolves from state $A$ to state $B$, it should
be possible to reverse all momenta at state $B$ and recover state $A$. This
appears to contradict the second law of thermodynamics, which asserts that
macroscopic processes are irreversible.

Conventional resolutions appeal to coarse-graining, the overwhelming probability
of high-entropy states, or special initial conditions at the Big Bang. We argue
that all such resolutions share a common flaw: they accept the premise of the
paradox --- that the velocity reversal operation is coherently definable --- and
then argue about probabilities. We instead show that the premise is incoherent,
and that the paradox dissolves before any probabilistic argument is needed.

\subsection{The Frozen Universe Requirement}
\label{subsec:frozen}

\begin{theorem}[Velocity Reversal Requires a Frozen Universe]
\label{thm:frozen}
The Loschmidt velocity reversal operation is coherently definable only in a
universe where the complement of the system under study is frozen --- that is,
does not progress --- during the reversal operation. A frozen universe is
impossible under the partition framework.
\end{theorem}

\begin{proof}
The velocity reversal operation consists of:
\begin{enumerate}
    \item Identifying the system $\mathcal{S}$ and its state $\{x_i(t_s),
          p_i(t_s)\}$ at reversal time $t_s$
    \item Constructing the reversed state $\{x_i(t_s), -p_i(t_s)\}$
    \item Releasing the system into this reversed state
\end{enumerate}
Step 1 requires observing the system, which is a partition event: $\mathcal{S}$
must be distinguished from $\overline{\mathcal{S}}$ to be observed. This
partition event takes time $\tau_{p,\mathrm{obs}} > 0$
(Postulate~\ref{post:duration}) and generates residue
$\beta_{\mathrm{obs}} > 0$ (Postulate~\ref{post:floor}).

During $\tau_{p,\mathrm{obs}}$, the complement $\overline{\mathcal{S}}$ must
progress (Postulate~\ref{post:duration}), advancing from state $U(t_s)$ to
$U(t_s + \tau_{p,\mathrm{obs}})$.

Step 2 requires computing the negated momenta, which takes additional time
$\tau_{p,\mathrm{compute}} > 0$ during which $\overline{\mathcal{S}}$ advances
to $U(t_s + \tau_{p,\mathrm{obs}} + \tau_{p,\mathrm{compute}})$.

Step 3 releases $\mathcal{S}$ into state $\{x_i(t_s), -p_i(t_s)\}$ --- a state
computed at $t_s$ --- into a universe now at
$U(t_s + \tau_{p,\mathrm{obs}} + \tau_{p,\mathrm{compute}})$, not $U(t_s)$. The
system is released into the wrong universe. The reversed state is only
time-reversed relative to $U(t_s)$; the universe has moved on to
$U(t_s + \Delta\tau)$ with $\Delta\tau > 0$.

For the reversal to be coherent, we would need $\overline{\mathcal{S}}$ to
remain at $U(t_s)$ throughout. But a frozen $\overline{\mathcal{S}}$ violates
Postulate~\ref{post:duration}: if $\overline{\mathcal{S}}$ does not progress,
no partition event within it can be well-defined, and the observation in Step 1
cannot be performed. The frozen universe is therefore impossible under the
partition framework, and the velocity reversal operation is incoherent.
\end{proof}

\subsection{The Swimmer Argument}
\label{subsec:swimmer}

We illustrate the abstract argument with a concrete case. Consider a swimmer
performing the backstroke in a pool. The system $\mathcal{S}$ is the
swimmer-pool system. The Loschmidt argument claims: reverse all velocities, and
the swimmer will reverse their trajectory back to the starting position.

\begin{theorem}[Causal Path Removal by Reversal]
\label{thm:swimmer}
Reversing all interactions in the swimmer-pool system does not time-reverse the
system; it removes the causal path that constituted the system's history.
\end{theorem}

\begin{proof}
The swimmer's stroke at time $t_0$ generates a pressure wave in the pool. This
pressure wave is a residue $\beta_{\mathrm{stroke}}$ deposited into the
pool-water contact chain. The pressure wave:
\begin{enumerate}
    \item Propagates to the pool walls at time $t_0 + L/c_{\mathrm{water}}$,
          generating residue $\beta_{\mathrm{wall}}$
    \item Reflects and propagates to the far side, generating residue
          $\beta_{\mathrm{reflect}}$
    \item Deposits thermal energy throughout the pool, generating residues
          $\{\beta_{\mathrm{thermal},k}\}$
    \item Propagates into the pool structure, the ground, and the
          surrounding environment, generating further residues
\end{enumerate}
Each residue $\beta_k$ is the causal propagator for the next event in the chain
(Theorem~\ref{thm:causal_residue}). The causal path of the stroke is the
sequence $\beta_{\mathrm{stroke}} \to \beta_{\mathrm{wall}} \to
\beta_{\mathrm{reflect}} \to \cdots$.

By time $t_1 > t_0$, the residues $\{\beta_k\}$ have propagated throughout the
contact chain of the universe. They are no longer localised in the pool. They
are in the pool structure, the ground, the air, the building, the Earth's crust.

To reverse the swimmer, one would need to reverse all momenta including those
of all items that have received any $\beta_k$. This includes the pool walls,
the pool structure, the ground, the air, and the Earth's crust. The reversal
requirement expands without bound through the contact chain.

Attempting to reverse a finite subset --- the pool and swimmer but not the
ground --- would be incomplete: the ground's contact chain now includes the
residues of the stroke and has advanced its partition count $M_{\mathrm{ground}}$
accordingly. The pool, once reversed, would be embedded in a ground that has
already received and propagated the stroke's residues. The contact chain between
pool and ground would generate forward-directed partition events immediately upon
contact, because the ground's state is not time-reversed.

Most fundamentally: the causal path of the stroke --- the sequence of residues
$\beta_k$ that propagated causally from the stroke --- cannot be reversed because
reversing it would require each residue to un-cause the next one. A residue
cannot un-cause a subsequent partition event: the partition event has already
established new contact relationships that themselves have residues. Un-causation
is not a physical operation; it would require the partition count $M$ to
decrease, which Theorem~\ref{thm:floor} forbids.

Therefore reversing all velocities does not time-reverse the system. It removes
the causal path that made the system a system.
\end{proof}

\subsection{The Spin Echo: A Forward Process Misread as Reversal}
\label{subsec:spin_echo}

The nuclear magnetic resonance spin echo appears to support time-reversal: after
a $\pi$ pulse at time $\tau$, the spins refocus at time $2\tau$ as if the
dephasing had been reversed. We show this is a forward process.

\begin{theorem}[Spin Echo as Forward Partition Process]
\label{thm:spin_echo}
The NMR spin echo is not a time-reversed process. It is a forward partition
process whose initial conditions were set by the residue of the $\pi$ pulse
interacting with the residue of the dephasing process.
\end{theorem}

\begin{proof}
The spin echo proceeds as follows in partition terms:
\begin{enumerate}
    \item \textbf{Excitation}: $\pi/2$ pulse establishes phase coherence.
          Partition count $M_0$.
    \item \textbf{Dephasing} (interval $[0,\tau]$): Each spin $k$ accumulates
          partition count at rate $\dot{M}_k = \omega_k/(2\pi)$ where $\omega_k
          = \omega_0 + \delta\omega_k$ includes local field inhomogeneity
          $\delta\omega_k$. At time $\tau$: partition count of spin $k$ is
          $M_k(\tau) = \dot{M}_k \cdot \tau$. This is a forward process.
    \item \textbf{$\pi$ pulse at $\tau$}: Inverts $z$-component of each spin.
          In partition terms: applies a parity transformation to the residue
          parity $s \to -s$ of each spin's partition state. This is a new
          partition event with residue $\beta_\pi > 0$. Partition count advances
          to $M(\tau) + 1$.
    \item \textbf{Refocusing} (interval $[\tau, 2\tau]$): Each spin continues
          accumulating partition count forward at the same rate $\dot{M}_k$.
          At time $2\tau$: partition count is $M_k(2\tau) = \dot{M}_k \cdot
          2\tau$. The phase of spin $k$ is:
          \begin{equation}
              \phi_k(2\tau) = \omega_k \cdot 2\tau - 2\omega_k\tau = 0
          \end{equation}
          (the $\pi$ pulse effectively negated the accumulated phase by flipping
          the partition parity, and the subsequent forward accumulation cancels
          the dephasing exactly).
\end{enumerate}
The refocusing is not time reversal. The partition count $M_k$ has increased
monotonically from $0$ to $\dot{M}_k \cdot 2\tau$ throughout the entire echo
sequence. Time has advanced by $2\tau$, not returned to zero. The coherence
restoration arises because the $\pi$ pulse is a new partition event whose
residue (the parity flip) is precisely structured to cancel the accumulated
phase --- not to reverse it.

The universe advanced its partition count by $\dot{M}_{\mathrm{universe}} \cdot
2\tau$ throughout the echo. The magnet, the sample tube, the laboratory, the
Earth, all progressed. The echo occurred inside a universe that had moved
forward by $2\tau$, not inside the universe at time 0.
\end{proof}

\subsection{The T2* and T2 Distinction as Partition Depth Split}
\label{subsec:T2_split}

The distinction between $T_2^*$ (apparent transverse relaxation time, including
field inhomogeneity) and $T_2$ (true transverse relaxation time) maps directly
onto the partition framework:

\begin{itemize}
    \item $T_2^*$ is the dephasing time at the \emph{microstate} level: each
          spin's individual partition trajectory diverges due to local field
          variations. These are reversible in the sense that the $\pi$ pulse
          can cancel them --- they are phase errors, not entropy, because the
          local field variations are themselves part of the partition chain and
          their residues are coherent.

    \item $T_2$ is the dephasing time at the \emph{partition count} level:
          irreversible loss of coherence due to spin-spin interactions that
          deposit residues $\beta_k$ into the contact chain between spins.
          These residues propagate into the broader contact chain and cannot
          be recovered by the $\pi$ pulse.
\end{itemize}

The ratio $T_2^*/T_2$ is the ratio of reversible phase accumulation (microstate
level, recoverable by parity flip) to irreversible partition count accumulation
(contact chain level, not recoverable). The spin echo recovers $T_2^*$; it
cannot recover $T_2$.

\subsection{Formal Resolution of the Paradox}
\label{subsec:formal_resolution}

\begin{theorem}[Loschmidt Paradox Resolution]
\label{thm:resolution}
There is no paradox between the time-reversal symmetry of microscopic equations
and the irreversibility of macroscopic processes. The apparent paradox arises
from a category error: the time-reversal symmetry of equations of motion applies
to the \emph{form} of the equations, not to the \emph{states} to which those
equations apply. The states are constrained by the partition framework to satisfy
$M(t_2) > M(t_1)$ for $t_2 > t_1$. No state trajectory with decreasing $M$ is
physically accessible, regardless of whether the equations of motion are
time-reversal symmetric.
\end{theorem}

\begin{proof}
Let $\mathcal{E}$ be the equations of motion of a system $\mathcal{S}$, and let
$T$ be the time-reversal operator: $T: (x(t), p(t)) \mapsto (x(-t), -p(-t))$.
The statement ``$\mathcal{E}$ is time-reversal symmetric'' means: if
$(x(t), p(t))$ satisfies $\mathcal{E}$, then $(x(-t), -p(-t))$ also satisfies
$\mathcal{E}$.

This is a statement about the \emph{solution space} of $\mathcal{E}$: both
$(x(t), p(t))$ and its time-reverse are mathematical solutions. It is not a
statement that the time-reversed solution is physically accessible from the
forward solution.

Physical accessibility requires:
\begin{enumerate}
    \item The initial conditions of the reversed solution must be achievable by
          a physical operation. This requires the velocity reversal operation,
          which Theorem~\ref{thm:frozen} shows is incoherent.
    \item The reversed solution must have $M_{\mathrm{reversed}}(t_2) >
          M_{\mathrm{reversed}}(t_1)$ for $t_2 > t_1$. A time-reversed solution
          $(x(-t), -p(-t))$ has decreasing coordinate time, which corresponds
          to decreasing partition count. This violates Theorem~\ref{thm:floor}.
\end{enumerate}
Therefore, while time-reversed solutions exist as mathematical objects
satisfying $\mathcal{E}$, no physical process can access them. The paradox
arises from confusing mathematical solutions with physical states. There is no
paradox about physical states: all physical states have monotone increasing $M$,
and no physical operation can violate this.
\end{proof}
